\documentclass[aspectratio=169]{beamer}

% ── Thème ────────────────────────────────────────────────────────────────
\usetheme{Singapore}
\setbeamercovered{invisible}

% ── Encodage et langue ────────────────────────────────────────────────────
\usepackage[utf8]{inputenc}
\usepackage[T1]{fontenc}
\usepackage[french]{babel}

% ── Packages ─────────────────────────────────────────────────────────────
\usepackage{amsmath, amssymb}
\usepackage{tikz}
\usetikzlibrary{arrows.meta, positioning}
\usepackage{xcolor}
\usepackage[normalem]{ulem}   % \sout pour le titre barré

% ── Couleurs dominante rouge ──────────────────────────────────────────────
\definecolor{mainred}{RGB}{165,15,15}
\definecolor{darkred}{RGB}{100,5,5}

\setbeamercolor{structure}{fg=mainred}
\setbeamercolor{frametitle}{fg=darkred, bg=red!10}
\setbeamercolor{block title}{bg=mainred, fg=white}
\setbeamercolor{block body}{bg=red!5, fg=black}
\setbeamercolor{alerted text}{fg=darkred}
\setbeamercolor{section in head/foot}{bg=darkred, fg=white}
\setbeamercolor{subsection in head/foot}{bg=mainred!70, fg=white}
\setbeamercolor{title}{fg=darkred}
\setbeamercolor{author}{fg=darkred!80}

% ── Métadonnées ───────────────────────────────────────────────────────────
\title[Parcours \& Thèse]{%
  Présentation de mon parcours\\[0.2em]
  et de mon projet de thèse}
\author{Anatole \textsc{Joseph-Stouls}}
\institute{ENS Rennes $\cdot$ LIED — Candidature EDPIF\\[0.2em]
  % BALISE_LOGO : logos ENS Rennes, LIED, CNRS, Université Paris Cité
}
\date{2025}

% ══════════════════════════════════════════════════════════════════════════
\begin{document}

% ── Titre ─────────────────────────────────────────────────────────────────
\begin{frame}
  \titlepage
\end{frame}

% ── Plan ──────────────────────────────────────────────────────────────────
\begin{frame}{Plan}
  \tableofcontents
\end{frame}

% ══════════════════════════════════════════════════════════════════════════
\section{Parcours}

% ─────────────────────────────────────────────────────────────────────────
% SLIDE 3 — Parcours
% TEXT: Très pluridisciplinaire (outil visuel : frise chronologique du CV)
% Classe prépa → ENSAM | j'aime la théorie, les maths du réel, je repasse
% en candidat libre | ENS (famille modeste, poste mécatronique accepté) |
% J'y ai notamment étudié : | Ingénierie Électrique, Mécanique, Informatique |
% traitement signal, politique énergétique, MMC, fluides, thermo (micro→macro) |
% À côté : licence et master maths fondamentales, proba-stats | chaînes de
% Markov, théorème ergodique | intérêts : économie hétérodoxe, sociologie |
% statistiques de l'INSEE | taxe carbone | stages : prof de maths prépa |
% CAO Munich | R&D tracking gaz fossile ||
% T1–T14 (14 transitions)
% ─────────────────────────────────────────────────────────────────────────
\begin{frame}[t,shrink=12]{Parcours}
\vspace{-0.4em}
% ── Frise chronologique ───────────────────────────────────────────────────
\resizebox{\linewidth}{!}{%
\begin{tikzpicture}
  \draw[-{Stealth[length=5pt]}, mainred, line width=1.2pt] (0,0) -- (12.8,0);
  % Blocs ères
  \fill[mainred!20, draw=mainred, thin] (0,0.1) rectangle (2.8,0.5);
  \node[font=\small, darkred] at (1.4,0.3) {\textbf{CPGE}};
  \fill[mainred!20, draw=mainred, thin] (2.8,0.1) rectangle (4.0,0.5);
  \node[font=\small, darkred] at (3.4,0.3) {ENSAM};
  \only<3->{\fill[mainred!40, draw=darkred, line width=0.8pt]
    (4.0,0.05) rectangle (8.4,0.6);}
  \only<3->{\node[font=\small\bfseries, darkred] at (6.2,0.32)
    {ENS Rennes};}
  \only<14->{\fill[mainred!20, draw=mainred, thin]
    (8.4,0.1) rectangle (10.0,0.5);}
  \only<14->{\node[font=\small, darkred] at (9.2,0.3) {ENSAM};}
  % Repères temporels
  \foreach \x/\yr in {0/2018,2.8/2020,4.0/2021,8.4/2024,10.0/2025}{
    \draw[mainred!50,thin] (\x,-0.08) -- (\x,0.08);
    \node[font=\tiny, mainred!70, below=1pt] at (\x,0) {\yr};
  }
\end{tikzpicture}}
\vspace{0.1em}
% ── Contenu ───────────────────────────────────────────────────────────────
\begin{columns}[t]
  \begin{column}{0.52\textwidth}
    \begin{itemize}[\footnotesize]\footnotesize\setlength\itemsep{1.5pt}
      \item<1-> CPGE MPSI puis PSI
        {\scriptsize(Lakanal / Michelet)}
        — concours $\to$ ENSAM Lille
      \item<2-> \textit{«\,j'aime la théorie, les maths du réel\,»}
        {\scriptsize$\to$ repasse les concours en candidat libre}
      \item<3-> \alert{\textbf{ENS Rennes}}, Dépt. Mécatronique
        {\scriptsize(élève fonctionnaire)}
      \item<4-> \textbf{Ingénierie :} élec., méca., info.
      \item<5-> {\scriptsize signal, énergie, MMC, fluides,
        thermo, électrotechnique — approche micro\,$\to$\,macro}
      \item<6-> \textbf{Maths :} L3 (Angers) + M1 fond. (Nantes)
        {\scriptsize— proba-stats}
      \item<7-> {\scriptsize chaînes de Markov, théorème ergodique —
        utiles dès aujourd'hui en physique statistique}
    \end{itemize}
  \end{column}
  \begin{column}{0.43\textwidth}
    \begin{itemize}[\footnotesize]\footnotesize\setlength\itemsep{1.5pt}
      \item<8-> \textbf{Centres d'intérêts :}
      \item<9-> {\scriptsize macro-éco hétérodoxe, finance,
        sociologie, géopolitique}
      \item<10-> {\scriptsize $\to$ comprendre les statistiques de l'INSEE}
      \item<11-> {\scriptsize $\to$ saisir la théorie derrière
        la taxe carbone}
      \item<12-> \textbf{Stages :}
      \item<13-> {\scriptsize Prof de maths en prépa
        \textit{(j'aime enseigner)}}
      \item<14-> {\scriptsize Ingénieur CAO, usinage —
        Cinetica GmbH, Munich}
      \item<15-> {\scriptsize R\&D tracking gaz fossile — Attributes}
    \end{itemize}
  \end{column}
\end{columns}
\end{frame}

% ══════════════════════════════════════════════════════════════════════════
\section{Recherche au LIED}

% ─────────────────────────────────────────────────────────────────────────
% SLIDE 4 — Stage LIED : l'effet rebond
% TEXT: le dernier stage, débuté en mars de cette année au | LIED |
% mission : expliquer l'effet rebond en | économie | à travers le prisme
% de la | physique | et de la | biologie : amélioration techno → baisse
% consommation... | mais effet rebond : décrue moindre, voire augmentation |
% (Hendrick) → hypothèse : phénomène endogène à une organisation | auto-critique
% | ville | être vivant | colonie de bactéries | tissu de dette |
% tous → mécanisme commun issu de l'organisation micro ||
% T15–T26 (12 transitions)
% ─────────────────────────────────────────────────────────────────────────
\begin{frame}[t]{Recherche au LIED — L'effet rebond}
\begin{columns}[T]
  \begin{column}{0.50\textwidth}
    \begin{block}{Le dernier stage\only<2->{ — \textbf{LIED}}}
      \footnotesize
      \only<1>{Stage débuté en mars de cette année.}
      \only<2->{LIED — UMR 8236, CNRS / Université Paris Cité.
        Encadrant : \textbf{Éric Herbert}.\\[0.2em]
        Mission : expliquer l'\alert{effet rebond}%
        \only<3->{ en \textbf{économie}}%
        \only<4->{, à travers le prisme de la \textbf{physique}}%
        \only<5->{ et de la \textbf{biologie}.}}
    \end{block}
    \vspace{0.3em}
    \begin{block}<6->{}
      \footnotesize
      Amélioration du rendement technologique\\
      $\to$ devrait faire \textbf{baisser} la consommation d'énergie.\\[0.2em]
      \uncover<7->{\alert{\textbf{Mais :}} la décrue est moindre qu'attendue.\\
      Parfois, la consommation \textbf{augmente} globalement.\\[0.1em]
      \scriptsize(Réf. Hendrick et al.)}
    \end{block}
  \end{column}
  \begin{column}{0.45\textwidth}
    \begin{block}<8->{Hypothèse}
      \footnotesize
      L'effet rebond serait endogène à une organisation
      \alert{\textbf{auto-critique}}\only<8>{ de la société.}
      \only<9->{.\\[0.3em]
      Comportements similaires entre :
      \begin{itemize}\setlength\itemsep{2pt}
        \item<10-> une \textbf{ville}
        \item<11-> un \textbf{être vivant}
        \item<12-> une \textbf{colonie de bactéries}
        \item<13-> un \textbf{tissu de dette scripturale}
      \end{itemize}
      \uncover<14->{\scriptsize Tous consomment énergie/matière,
        produisent déchets et information.\\
        $\to$ mécanisme commun issu de l'organisation \textbf{micro}.}}
    \end{block}
  \end{column}
\end{columns}
\end{frame}

% ─────────────────────────────────────────────────────────────────────────
% SLIDE 5 — Système auto-critique : l'avalanche
% TEXT: Un système auto-critique est... | Exemple : grille |
% flocons se déposent | - | - | - | colonne → seuil → topple |
% si voisin atteint seuil : avalanche taille 2 | état critique atteint |
% des petites avalanches | des grosses | distribution = loi de puissance |
% règle simple, comportement complexe ||
% T27–T39 (13 transitions)
% ─────────────────────────────────────────────────────────────────────────
\begin{frame}[t]{Système auto-critique — L'avalanche de sable}
\begin{columns}[T]
  \begin{column}{0.50\textwidth}
    \only<1>{%
      \begin{block}{Définition}
        \small
        Un système dynamique qui évolue
        \textbf{spontanément} vers un état critique
        où de petites perturbations provoquent
        des événements de \textbf{toutes tailles},
        distribués en \textbf{loi de puissance} —
        indépendamment des paramètres exogènes.
      \end{block}}
    \only<2-12>{%
      \vspace{0.3em}\centering
      % ── Sandpile TikZ (5 colonnes, seuil = 4) ───────────────────────
      % Hauteurs : états successifs
      % T28 grille vide, T29 [0,0,1,0,0], T30 [1,0,1,0,0]
      % T31 [1,0,2,0,0], T32 [1,0,3,0,0]
      % T33 col3 à 4 (rouge), T34 après topple [1,1,2,1,0]
      % T35 état dense, T36 petite aval, T37 grande cascade
      \begin{tikzpicture}[scale=0.9]
        % Fond de grille
        \foreach \c in {0,...,4}{
          \foreach \r in {0,...,3}{
            \draw[mainred!15, thin]
              ({\c*0.85},{(\r)*0.52}) rectangle ++(0.8,0.47);}}
        % Seuil (=4)
        \draw[darkred, dashed, thin]
          (-0.1,{3*0.52+0.5}) -- ({4.6},{3*0.52+0.5})
          node[right,font=\tiny,darkred]{seuil\,4};
        % Grains par état (\only<n>{grain})
        % Grains persistants :
        % T29 : col 3 (idx 2) rang 0
        \only<3->{\fill[mainred!65]({2*0.85},0) rectangle ++(0.8,0.47);}
        % T30 : col 1 (idx 0) rang 0
        \only<4->{\fill[mainred!65](0,0) rectangle ++(0.8,0.47);}
        % T31 : col 3 rang 1
        \only<5->{\fill[mainred!65]({2*0.85},0.52) rectangle ++(0.8,0.47);}
        % T32 : col 3 rang 2
        \only<6->{\fill[mainred!65]({2*0.85},{2*0.52}) rectangle ++(0.8,0.47);}
        % T33 : col 3 rang 3 (critique, rouge vif)
        \only<7>{\fill[red!90]({2*0.85},{3*0.52}) rectangle ++(0.8,0.47);
          \node[font=\tiny\bfseries,red] at ({2*0.85+0.4},{4*0.52})
            {topple !};}
        % T34 : après topple → [1,1,2,1,0]
        \only<8->{
          % col 2 (idx 1) rang 0
          \fill[mainred!50](0.85,0) rectangle ++(0.8,0.47);
          % col 3 rang 0 et 1 (reste)
          \fill[mainred!65]({2*0.85},0) rectangle ++(0.8,0.47);
          \fill[mainred!65]({2*0.85},0.52) rectangle ++(0.8,0.47);
          % col 4 (idx 3) rang 0
          \fill[mainred!50]({3*0.85},0) rectangle ++(0.8,0.47);}
        % T35 : état critique (beaucoup de grains)
        \only<9->{
          \fill[mainred!60](0,0.52) rectangle ++(0.8,0.47);
          \fill[mainred!60](0.85,0.52) rectangle ++(0.8,0.47);
          \fill[mainred!60](0.85,{2*0.52}) rectangle ++(0.8,0.47);
          \fill[mainred!60]({2*0.85},{2*0.52}) rectangle ++(0.8,0.47);
          \fill[mainred!60]({3*0.85},0.52) rectangle ++(0.8,0.47);
          \fill[mainred!60]({3*0.85},{2*0.52}) rectangle ++(0.8,0.47);
          \fill[mainred!60]({4*0.85},0) rectangle ++(0.8,0.47);}
        % T36 : petite avalanche (1 topple)
        \only<10>{\fill[red!80]({1*0.85},{3*0.52}) rectangle ++(0.8,0.47);
          \node[font=\tiny,red] at ({0.85+0.4},{4*0.52+0.05}){topple};}
        % T37 : grande cascade (3 topples)
        \only<11>{
          \fill[red!80]({1*0.85},{3*0.52}) rectangle ++(0.8,0.47);
          \fill[red!60]({2*0.85},{3*0.52}) rectangle ++(0.8,0.47);
          \fill[red!60](0,{3*0.52}) rectangle ++(0.8,0.47);
          \draw[->,red,thick]({0.85+0.4},{3*0.52+0.48})--
            ({0.4},{3*0.52+0.48});
          \draw[->,red,thick]({0.85+0.4},{3*0.52+0.48})--
            ({2*0.85+0.4},{3*0.52+0.48});}
        % Labels colonnes
        \foreach \c in {0,...,4}{
          \node[font=\tiny,mainred!80] at ({\c*0.85+0.4},-0.25){\c};}
      \end{tikzpicture}}
  \end{column}
  \begin{column}{0.45\textwidth}
    \footnotesize
    \begin{itemize}\setlength\itemsep{5pt}
      \item<10-> Petites avalanches : \textbf{fréquentes}
      \item<11-> Grandes avalanches : \textbf{rares}
      \item<12-> Distribution statistique :\\[0.2em]
        \alert{loi de puissance} $P(s)\sim s^{-\alpha}$\\[0.2em]
        % BALISE_GRAPH_01 : graphe fréquence vs taille (log-log)
        {\scriptsize\textit{[BALISE\_GRAPH\_01]}}
    \end{itemize}
    \vspace{0.3em}
    \begin{alertblock}<13->{}
      \footnotesize\centering
      \textbf{Règle simple $\to$ comportement complexe.}\\[0.2em]
      Insensibilité aux paramètres exogènes.
    \end{alertblock}
  \end{column}
\end{columns}
\end{frame}

% ─────────────────────────────────────────────────────────────────────────
% SLIDE 6 — La simulation
% TEXT: J'ai codé une simulation | entités + environnement |
% entités apparaissent | - | - | extraient énergie de l'environnement |
% fonction concave de leur | taille actuelle ; énergie se |
% dissipe exponentiellement | peuvent se prêter de l'énergie |
% si rentes > ressources, une entité | n'a plus d'énergie | disparaît |
% résultats : n_alive stable + 2 populations | formaliser ? ||
% T40–T54 (15 transitions)
% ─────────────────────────────────────────────────────────────────────────
\begin{frame}[t]{La simulation}
\begin{columns}[T]
  \begin{column}{0.48\textwidth}
    \begin{block}{Règles du modèle}
      \footnotesize
      \begin{itemize}\setlength\itemsep{2pt}
        \item<1-> Règles simples, naturelles (rasoir d'Occam)
        \item<2-> \textbf{Entités} interagissent entre elles
          et avec leur \textbf{environnement}
        \item<3-> Entités \textbf{apparaissent} continuellement
          ($\lambda$ constant)
        \item<6-> Extraient de l'énergie de l'environnement
        \item<7-> Puissance extractrice = $f$\,\textbf{concave}(taille)
        \item<9-> Énergie se \textbf{dissipe} exponentiellement
        \item<10-> Peuvent se \textbf{prêter} de l'énergie
          contre une rente
        \item<12-> Insolvabilité \only<13->{$\to$}
          \only<14->{\alert{\textbf{disparition}}}
      \end{itemize}
    \end{block}
  \end{column}
  \begin{column}{0.47\textwidth}
    % ── Visualisation réseau d'entités ───────────────────────────────
    \only<1-13>{%
      \centering\vspace{0.4em}
      \begin{tikzpicture}[scale=0.85,
        ent/.style={circle,draw=mainred,fill=red!15,
          minimum size=0.6cm,font=\tiny\bfseries},
        dead/.style={circle,draw=gray!60,fill=gray!10,
          minimum size=0.6cm,font=\tiny,dashed}]
        % Environnement
        \draw[mainred!25,fill=blue!4,rounded corners=4pt]
          (-2.2,-2.0) rectangle (2.2,2.0);
        \node[font=\scriptsize,mainred!60] at (0,-1.75)
          {Environnement};
        % Entités apparaissent progressivement
        \node[ent](e1) at (0.1,0.9) {$e_1$};
        \only<4->{\node[ent](e2) at (-1.1,-0.2) {$e_2$};}
        \only<5->{\node[ent](e3) at (1.1,-0.3) {$e_3$};}
        % Extraction d'énergie
        \only<6->{\draw[->,mainred!70,thick]
          (-1.9,1.0)--(e1.west)
          node[midway,above,font=\tiny]{$\Pi$};}
        % Prêt entre entités
        \only<10->{\draw[->,darkred,dashed,bend right=20]
          (e1) to node[right,font=\tiny]{prêt} (e3);}
        % Entité en danger
        \only<11>{\draw[red,thick] (e3.south west)--(e3.north east);
          \draw[red,thick] (e3.north west)--(e3.south east);}
        % Entité disparue
        \only<12->{\node[dead] at (1.1,-0.3) {$\times$};}
      \end{tikzpicture}}
    % ── Résultats ─────────────────────────────────────────────────────
    \only<14-15>{%
      \begin{block}{Résultats}
        \footnotesize
        \begin{itemize}\setlength\itemsep{4pt}
          \item $n_\text{alive}$ \textbf{ne diverge pas}
            malgré le flux de naissance constant
          \item Émergence de \textbf{deux populations}
            aux comportements économiques distincts
        \end{itemize}
        \vspace{0.3em}
        % BALISE_GRAPH_02 : n_alive(t) et distribution Gini des actifs
        {\scriptsize\textit{[BALISE\_GRAPH\_02 — rapport\_final\_sensibilite]}}
      \end{block}
      \uncover<15->{%
        \begin{alertblock}{}
          \small\centering
          Comment expliquer les macro\\
          à partir des règles micro ?
        \end{alertblock}}}
  \end{column}
\end{columns}
\end{frame}

% ─────────────────────────────────────────────────────────────────────────
% SLIDE 7 — Physique statistique
% TEXT: On a donc un système | graphe orienté pondéré |
% myriade de micro-états | qui évoluent → macro-états |
% je suis en plein dans la physique statistique. ||
% T55–T59 (5 transitions)
% ─────────────────────────────────────────────────────────────────────────
\begin{frame}{Physique statistique}
\begin{center}
\begin{tikzpicture}[scale=0.88,
  nd/.style={circle,draw=mainred,fill=red!12,
    minimum size=0.75cm,font=\scriptsize}]
  % Boîte système
  \uncover<1->{\draw[mainred,rounded corners,thick]
    (-3.8,-1.9) rectangle (3.6,2.0);
    \node[font=\small,darkred,above] at (0,2.0) {\textbf{Système}};}
  % Graphe orienté pondéré
  \uncover<2->{
    \node[nd](n1) at (-2.4,1.0) {$i$};
    \node[nd](n2) at (-0.5,1.3) {$j$};
    \node[nd](n3) at (1.5,0.9) {$k$};
    \node[nd](n4) at (-1.5,-0.8) {$l$};
    \node[nd](n5) at (1.0,-1.1) {$m$};
    \draw[->,mainred,thick](n1)--(n2)
      node[midway,above,font=\tiny]{$w_{ij}$};
    \draw[->,mainred,thick](n2)--(n3);
    \draw[->,mainred!70](n3)--(n5);
    \draw[->,mainred!70](n4)--(n1);
    \draw[->,mainred!70](n5)--(n4);
    \node[font=\scriptsize,darkred] at (0,-1.7)
      {graphe orienté pondéré};}
  % Micro-états
  \uncover<3->{
    \node[font=\scriptsize,mainred!80] at (-3.0,-1.5)
      {$|\Omega_\text{micro}| \gg 1$};}
  % Macro flèche + macro-états
  \uncover<4->{
    \draw[->,darkred,line width=2.5pt](3.8,0.05)--(5.4,0.05);
    \node[draw=darkred,fill=red!15,rounded corners,
      font=\small\bfseries,darkred,align=center,
      minimum width=2.4cm] at (6.8,0.05)
      {macro-états\\{\scriptsize distributions}};}
\end{tikzpicture}
\end{center}
\vspace{-0.3em}
\uncover<5->{%
  \begin{center}
    {\Large\color{mainred}\textbf{Physique statistique}}
  \end{center}}
\end{frame}

% ══════════════════════════════════════════════════════════════════════════
\section{Sujet de thèse}

% ─────────────────────────────────────────────────────────────────────────
% SLIDE 8 — État de l'art
% TEXT: des gens ont déjà réfléchi à la question : Oui, |
% Yakovenko APS Rev Mod Phys | Boltzmann-Gibbs + Pareto |
% En France, | Bouchaud & Mézard, loi de puissance |
% Institut Santa Fe | thèse Polytechnique 2022 |
% "Grandes fluctuations macroéconomiques..." ||
% T60–T67 (8 transitions)
% ─────────────────────────────────────────────────────────────────────────
\begin{frame}[t]{État de l'art}
\begin{columns}[T]
  \begin{column}{0.52\textwidth}
    \footnotesize
    Des gens ont déjà réfléchi à la question.
    \begin{itemize}\setlength\itemsep{5pt}
      \item<2-> \textbf{Yakovenko \& Rosser},
        \emph{Rev.\ Mod.\ Phys.} \textbf{81} (2009) :\\
        \uncover<3->{distributions de capital $=$ \textbf{Boltzmann-Gibbs}
          (bas revenus) + \textbf{Pareto} (hauts revenus)\\
          % BALISE_GRAPH_03 : figure Yakovenko Rev Mod Phys
          {\scriptsize\textit{[BALISE\_GRAPH\_03]}}}
      \item<4-> \textbf{Bouchaud \& Mézard},
        \emph{Physica A} \textbf{282} (2000) :\\
        \uncover<5->{distribution de capital = loi de puissance\\
          % BALISE_GRAPH_04 : figure Bouchaud Mézard
          {\scriptsize\textit{[BALISE\_GRAPH\_04]}}}
      \item<6-> \textbf{Institut Santa Fe} —
        physique des systèmes complexes économiques
    \end{itemize}
  \end{column}
  \begin{column}{0.43\textwidth}
    \begin{block}<7->{Thèse récente — Polytechnique, 2022}
      \footnotesize
      \uncover<8->{\og~\textbf{Grandes fluctuations
        macroéconomiques : criticité auto-organisée
        dans les réseaux d'entreprises, modèles à agents
        et matrices aléatoires}~\fg}
    \end{block}
    \vspace{0.4em}
    \begin{alertblock}<8->{}
      \footnotesize\centering
      C'est porteur. C'est d'actualité.
    \end{alertblock}
  \end{column}
\end{columns}
\end{frame}

% ─────────────────────────────────────────────────────────────────────────
% SLIDE 9 — Sujet de thèse
% TEXT: questions se poursuivent en thèse, m'intéresser à |
% organisation et dissipation dans une structure économique. Pas restreint
% à la physico-économie : polymères orientés | modèle d'Ising | ... |
% on arrive au titre : | Structures dissipatives hors équilibre... |
% deux montagnes : forte interaction | hors équilibre ||
% T68–T75 (8 transitions)
% BALISE_TITRE : à la transition T72, on pourra barrer un ancien titre
%   et le remplacer par le titre définitif (commande \sout disponible)
% ─────────────────────────────────────────────────────────────────────────
\begin{frame}[t]{Sujet de thèse}
\begin{columns}[T]
  \begin{column}{0.50\textwidth}
    \begin{block}<1->{Questions qui se poursuivent en thèse}
      \footnotesize
      \uncover<1->{Organisation et dissipation dans
        une structure économique.}\\[0.2em]
      \uncover<2->{Mais \textbf{pas restreint} à la physico-économie :}
      \begin{itemize}\setlength\itemsep{1pt}
        \item<3-> physique de la matière —
          \textbf{polymères orientés}
        \item<4-> \textbf{modèle d'Ising}
        \item<5-> \dots
      \end{itemize}
      \uncover<5->{\scriptsize
        Entités qui interagissent localement → même formalisme.}
    \end{block}
    \vspace{0.3em}
    \only<6>{%
      \begin{block}{}
        \footnotesize\centering
        On arrive au titre de la thèse\dots
      \end{block}}
    \only<7->{%
      \begin{block}{}
        \footnotesize\centering
        % BALISE_TITRE : remplacer/barrer ce titre si nécessaire (T72)
        \textbf{Structures dissipatives hors-équilibre,\\
          distributions stationnaires\\
          et flux d'énergie-matière}
      \end{block}}
  \end{column}
  \begin{column}{0.45\textwidth}
    \begin{block}<8->{Deux montagnes à gravir}
      \footnotesize
      \begin{enumerate}\setlength\itemsep{6pt}
        \item<8-> \textbf{Interactions fortes} entre entités\\
          {\scriptsize(vs. interactions faibles en méca-stats classique)}
        \item<9-> Structures \textbf{hors-équilibre}\\
          {\scriptsize(sujet fécond dans de nombreux labos
          de physique statistique)}
      \end{enumerate}
    \end{block}
  \end{column}
\end{columns}
\end{frame}

% ─────────────────────────────────────────────────────────────────────────
% SLIDE 10 — Objectifs et équipe
% TEXT: J'ai deux | objectifs | 1. systèmes auto-critiques → théorèmes |
% 2. modéliser société économique comme structure dissipative |
% continuité travaux LIED | Herbert + Giraud |
% ENS + proba maths | Goupil + Bruneton ||
% T76–T82 (7 transitions)
% ─────────────────────────────────────────────────────────────────────────
\begin{frame}[t]{Objectifs et équipe}
\begin{columns}[T]
  \begin{column}{0.50\textwidth}
    \begin{block}{Deux\uncover<2->{ objectifs}}
      \footnotesize
      \begin{enumerate}\setlength\itemsep{5pt}
        \item<3-> Étude formelle des systèmes auto-critiques
          — trouver des \textbf{théorèmes de physique statistique}
          {\scriptsize(potentiellement de niche, mais ce serait super)}
        \item<4-> Modéliser la société économique en tant que
          \textbf{structure dissipative}
          — données réelles + outils théoriques
      \end{enumerate}
    \end{block}
    \vspace{0.3em}
    \begin{block}<5->{Continuité des travaux du LIED}
      \footnotesize
      \uncover<6->{\textbf{Éric Herbert} (physique) +
        \textbf{Gaël Giraud} (économie) —
        modèle novateur de thermodynamique économique.}
    \end{block}
  \end{column}
  \begin{column}{0.45\textwidth}
    \begin{block}<7->{Atouts de mon parcours}
      \footnotesize
      Formation \textbf{concrète} ENS (ingénierie)
      $+$ \textbf{proba dures} (maths)
      $=$ vision holistique, théorique
      et phénoménologique.
    \end{block}
    \vspace{0.3em}
    \begin{block}<8->{Encadrement}
      \footnotesize
      \textbf{Christophe Goupil} et
      \textbf{Jean-Philippe Bruneton} —\\[0.2em]
      {\scriptsize ne laisseront passer, je pense,
      aucune faiblesse d'esprit.}
    \end{block}
  \end{column}
\end{columns}
\end{frame}

% ─────────────────────────────────────────────────────────────────────────
% SLIDE 11 — Un projet avec un début, un milieu et une fin
% TEXT: ce projet a un début, un milieu, et une fin. |
% Le début : travaux existants — Yakovenko, Bouchaud, Santa Fe |
% Le milieu : cette recherche |
% La fin : modèle robuste + outils théoriques + biblio ||
% T83–T85 (3 transitions)
% ─────────────────────────────────────────────────────────────────────────
\begin{frame}{Un projet structuré}
\begin{block}{}
  \centering\small
  Ce projet a un \textbf{début}, un \textbf{milieu}, et une \textbf{fin}.
\end{block}
\vspace{0.4em}
\begin{columns}[T]
  \begin{column}{0.30\textwidth}
    \begin{block}<2->{\centering Début}
      \footnotesize\centering
      Travaux existants\\[0.3em]
      Yakovenko, Bouchaud,\\
      Institut Santa Fe,\\
      thèse Polytechnique 2022
    \end{block}
  \end{column}
  \begin{column}{0.30\textwidth}
    \begin{block}<3->{\centering Milieu}
      \footnotesize\centering
      \textbf{Cette recherche}\\[0.3em]
      Structures dissipatives\\
      hors-équilibre,\\
      distributions stationnaires,\\
      flux énergie-matière
    \end{block}
  \end{column}
  \begin{column}{0.30\textwidth}
    \begin{block}<4->{\centering Fin}
      \footnotesize\centering
      Modèle physico-économique\\
      robuste\\[0.3em]
      Données réelles +\\
      outils théoriques +\\
      bibliographie
    \end{block}
  \end{column}
\end{columns}
\end{frame}

% ══════════════════════════════════════════════════════════════════════════
\section{Conclusion}

% ─────────────────────────────────────────────────────────────────────────
% SLIDE 12 — Après la thèse
% TEXT: MCF à l'université | agrégation de mathématiques, oraux
% dans quelques semaines, thèse requise pour enseigner en prépa ||
% T86–T87 (2 transitions)
% ─────────────────────────────────────────────────────────────────────────
\begin{frame}{Après la thèse}
\begin{columns}[T]
  \begin{column}{0.46\textwidth}
    \begin{block}<1->{Idéalement}
      \small
      Poste de \textbf{MCF à l'université} —
      si je n'ai pas à choisir entre
      recherche et enseignement.\\[0.4em]
      {\footnotesize J'aime vraiment enseigner.}
    \end{block}
  \end{column}
  \begin{column}{0.46\textwidth}
    \begin{block}<2->{En parallèle}
      \small
      \textbf{Agrégation de mathématiques}\\[0.3em]
      {\footnotesize Oraux dans quelques semaines.\\
      Une thèse est requise pour enseigner
      en classe préparatoire.}
    \end{block}
  \end{column}
\end{columns}
\end{frame}

% ─────────────────────────────────────────────────────────────────────────
% SLIDE 13 — Merci
% ─────────────────────────────────────────────────────────────────────────
\begin{frame}[plain]
\begin{center}
  \vspace{2.5em}
  {\Large\color{mainred}\textbf{Je vous remercie de votre attention}}\\[1.5em]
  {\large Questions ?}\\[1.5em]
  \textcolor{mainred!50}{\rule{0.55\textwidth}{0.5pt}}\\[0.8em]
  \small
  \textbf{Anatole Joseph-Stouls}\\[0.3em]
  \texttt{anatole.joseph-stouls@ens-rennes.fr}
\end{center}
\end{frame}

% ══════════════════════════════════════════════════════════════════════════
% ANNEXES
% ══════════════════════════════════════════════════════════════════════════
\appendix

% ─────────────────────────────────────────────────────────────────────────
% SLIDE 14 — Bibliographie
% ─────────────────────────────────────────────────────────────────────────
\begin{frame}{Bibliographie}
\footnotesize
\begin{enumerate}\setlength\itemsep{4pt}
  \item \textbf{Yakovenko \& Rosser} — Statistical mechanics of money,
    wealth, and income. \emph{Rev.\ Mod.\ Phys.} \textbf{81}, 1703 (2009).
  \item \textbf{Drăgulescu \& Yakovenko} — Statistical mechanics of money.
    \emph{Eur.\ Phys.\ J.\ B} \textbf{17}, 723 (2000).
  \item \textbf{Bouchaud \& Mézard} — Wealth condensation in a simple
    model of economy. \emph{Physica A} \textbf{282}, 536 (2000).
  \item \textbf{Herbert, Giraud, Louis-Napoléon, Goupil} — Macroeconomic
    dynamics in a finite world based on thermodynamic potential.
    \emph{Sci.\ Reports} \textbf{13}, 18020 (2023).
  \item \textbf{Touchette} — The large deviation approach to statistical
    mechanics. \emph{Phys.\ Rep.} \textbf{478}, 1 (2009).
  \item Thèse Polytechnique, 2022 — \emph{Grandes fluctuations
    macroéconomiques : criticité auto-organisée dans les réseaux
    d'entreprises, modèles à agents et matrices aléatoires.}
\end{enumerate}
\end{frame}

% ─────────────────────────────────────────────────────────────────────────
% SLIDE 15 — Annexe ENS
% ─────────────────────────────────────────────────────────────────────────
\begin{frame}{Annexe — Formation ENS Rennes (cours)}
\begin{columns}[T]
  \begin{column}{0.46\textwidth}
    \begin{block}{Ingénierie \& Physique}
      \footnotesize
      Électronique de puissance, conversion d'énergie,
      asservissement, mécanique des fluides,
      thermodynamique et phénomènes de transport,
      modélisation multi-physiques,
      élasto-visco-plasticité, dynamique des structures
    \end{block}
    \vspace{0.3em}
    \begin{block}{Informatique}
      \footnotesize
      Algorithmique, outils numériques,
      traitement d'images, vision par ordinateur,
      développement logiciel objet
    \end{block}
  \end{column}
  \begin{column}{0.46\textwidth}
    \begin{block}{Mathématiques}
      \footnotesize
      Analyse, algèbre, probabilités {\scriptsize(L3)} ;
      EDP, processus stochastiques {\scriptsize(M1)} ;
      théorie des grandes déviations,
      analyse fonctionnelle {\scriptsize(M2)}
    \end{block}
    \vspace{0.3em}
    \begin{block}{Transversal}
      \footnotesize
      Didactique, écoconception, marchés de l'énergie,
      robotique, mécanique analytique, RDM
    \end{block}
  \end{column}
\end{columns}
\end{frame}

% ─────────────────────────────────────────────────────────────────────────
% SLIDE 16 — Annexe Stage
% ─────────────────────────────────────────────────────────────────────────
\begin{frame}{Annexe — Résultats de la simulation}
\begin{columns}[T]
  \begin{column}{0.46\textwidth}
    \begin{block}{$n_\text{alive}$ ne diverge pas}
      \footnotesize
      Malgré un taux de création $\lambda$ constant,
      le nombre d'entités vivantes se stabilise.\\[0.3em]
      % BALISE_GRAPH_05 : n_alive(t) pour différentes valeurs de k
      {\scriptsize\textit{[BALISE\_GRAPH\_05]}}
    \end{block}
  \end{column}
  \begin{column}{0.46\textwidth}
    \begin{block}{Deux populations émergentes}
      \footnotesize
      Émergence spontanée de deux classes
      aux comportements économiques distincts.\\[0.3em]
      % BALISE_GRAPH_06 : distribution des actifs / Gini
      {\scriptsize\textit{[BALISE\_GRAPH\_06]}}
    \end{block}
  \end{column}
\end{columns}
\vspace{0.4em}
\begin{block}{Paramètres principaux}
  \footnotesize
  $\lambda$ (taux de création),
  $\theta$ (part du gain extrait),
  $k$ (candidats par emprunteur),
  $\mu$ (prime de rendement),
  $\delta$ (dépréciation),
  $\sigma_\alpha$ (volatilité brownienne)
\end{block}
\end{frame}

\end{document}
